\section{From the mean-square estimate to the zero-free region}\label{sec:reduction}

We first carry out Steps~1 and~2 of the outline: extract
cancellation in a M\"obius sum from a mean-square estimate for the family,
then use a Mellin transform to deduce nonvanishing.
Fix a Dirichlet character $\chi$ of modulus $q$, put $\nu=\chi\circ\NK$, and fix a finite set $S$
of prime ideals that contains all prime ideals above $2$ or $3$, as
well as all prime ideals dividing the conductor of $\nu$.
(In the formalization, $S$ consists of the prime ideals dividing $2592q$, where $2592=2^5\cdot3^4$.)
For an ideal or element $a$, write $(a,S)=1$ if no prime ideal in $S$
divides $a$.
An ideal is supported on $S$ if all its prime factors belong to $S$.
For $W\in C_c^\infty((0,\infty);\mathbb C)$, recall the family
\[
 A_u(D):=\sum_{(n,S)=1}\mu(n)\nu(n)\chi_n(u)W(\NK(n)/D),
\]
introduced in \eqref{eq:intro-family}, and the weighted M\"obius sum
\begin{equation}\label{eq:A-phi}
 A^\varphi(D):=\sum_{(n,S)=1}\mu(n)\nu(n)\frac{\varphi(n)}{\NK(n)}W(\NK(n)/D),
 \qquad
 \frac{\varphi(n)}{\NK(n)}=\prod_{\mathfrak p\mid n}\Bigl(1-\frac1{\NK(\mathfrak p)}\Bigr).
\end{equation}
Here and throughout the
original family, $n$ runs over ideals prime to $S$, represented by
their primary generators. The key estimate is the following.

\begin{proposition}\label{thm:ms}
Let $\vartheta=1/10$. For every fixed $\varepsilon>0$, there exists an integer $k=k(\varepsilon)\ge1$
such that, for every compact interval $I\subset(0,\infty)$,
all smooth $W$ supported in $I$, and $D\ge2$,
\begin{equation}\label{eq:ms}
 \sum_{0<\NK(u)\le D^{1+\vartheta}}|A_u(D)|^2
       \ll_{\nu,S,I,\varepsilon}
       \Bigl(\max_{0\le j\le k}\|W^{(j)}\|_\infty\Bigr)^2
       D^{2+\vartheta+\varepsilon}.
\end{equation}
\end{proposition}

We first deduce Theorem~\ref{thm:main} assuming Proposition~\ref{thm:ms}.
The proof of the proposition is completed in Section~\ref{sec:completion}.
The deduction uses \eqref{eq:ms} only for one $W$ at a time, with an implied constant that may
depend on $W$; this is the form that the formalization proves.

\subsection{Extraction}
\begin{proposition}\label{prop:extraction}
Let $I\subset(0,\infty)$ be a compact interval, let $W$ be smooth and supported in $I$, and let
$\varepsilon>0$. Suppose that \eqref{eq:ms} holds for $W$ with $2\varepsilon$ in place of
$\varepsilon$. Then, for $D\ge2$,
\begin{equation}\label{eq:mobius-saving}
 A^\varphi(D)\ll_{\nu,S,W,\varepsilon}D^{23/24+\varepsilon}.
\end{equation}
\end{proposition}

\begin{proof}
Put $H=D^{1+\vartheta}$ and
$Y=H^{1/6}=D^{(1+\vartheta)/6}=D^{11/60}$. For $b\in\OO$ and a prime
$\mathfrak p\notin S$, the character $\chi_{\mathfrak p}$ has order six, so
\begin{equation}\label{eq:sixth-power-rows}
 \chi_n(b^6)=\ind_{(b,n)=1}.
\end{equation}
Fix a smooth $V:\R\to\C$ supported in $[1/2,1]$, with $|V|\le1$ and
\[
 \widehat V(0):=\frac2{\sqrt3}\int_{\C}V(|z|^2)\dd x\dd y\ne0,
\]
the value at $0$ of the radial Fourier transform of Lemma~\ref{lem:poisson}, and consider
\[
 \mathcal V(D):=\sum_{b\in\OO}V\Bigl(\frac{\NK(b)}Y\Bigr)A_{b^6}(D)
 =\sum_{(n,S)=1}\mu(n)\nu(n)W(\NK(n)/D)
 \sum_{\substack{b\in\OO\\(b,n)=1}}V\Bigl(\frac{\NK(b)}Y\Bigr),
\]
where the second expression follows from \eqref{eq:sixth-power-rows}.
For squarefree $n$, inclusion--exclusion over $\mathfrak d\mid n$ writes the
inner sum as
\[
 \sum_{\mathfrak d\mid n}\mu(\mathfrak d)\sum_{c\in\OO}V\bigl(\NK(\mathfrak d)\NK(c)/Y\bigr).
\]
By Poisson summation (Lemma~\ref{lem:poisson}), the sum over $c$ is
$(Y/\NK(\mathfrak d))\widehat V(0)+O_V(1)$: the zero frequency gives the main term, the other
frequencies are $O_V(1)$ by the rapid decay of $\widehat V$, and when $\NK(\mathfrak d)>Y$ only
$c=0$ is in the support, where $V(0)=0$. Summing over $\mathfrak d$, and using
$2^{\omega(n)}\ll_\varepsilon\NK(n)^\varepsilon$ and $\NK(n)\ll_W D$, gives
\begin{equation}\label{eq:extraction-evaluation}
 \mathcal V(D)=Y\widehat V(0)A^\varphi(D)+O_{W,V,\varepsilon}(D^{1+\varepsilon}).
\end{equation}
The $b$ in the support of $V(\NK(b)/Y)$ satisfy $Y/2\le\NK(b)\le Y$; there are $O(Y)$ of them,
and their sixth powers are rows of norm at most $H$, each taken at most six times. Apply
\eqref{eq:ms} with loss $2\varepsilon$ and Cauchy--Schwarz to obtain
\begin{align}
 |A^\varphi(D)|^2
 &\ll \frac1{Y^2}\Bigl(Y\sum_{0<\NK(u)\le H}|A_u(D)|^2+D^{2+2\varepsilon}\Bigr)\nonumber\\
 &\ll_{\nu,S,W,\varepsilon}D^{2+\vartheta+2\varepsilon}/Y+D^{2+2\varepsilon}/Y^2
   =D^{23/12+2\varepsilon}+D^{49/30+2\varepsilon}.
                                                        \label{eq:prime-extract}
\end{align}
Thus $A^\varphi(D)\ll_{\nu,S,W,\varepsilon}D^{23/24+\varepsilon}$.
\end{proof}

\subsection{The Mellin transform}
For $\Rea s>0$ put
\[
 E_\varphi(s,\nu):=\prod_{\mathfrak p\notin S}
 \frac{1-\nu(\mathfrak p)\NK(\mathfrak p)^{-s}}{1-\nu(\mathfrak p)\NK(\mathfrak p)^{-s}(1-\NK(\mathfrak p)^{-1})},
 \qquad
 P_S(s,\nu):=\prod_{\mathfrak p\in S}\bigl(1-\nu(\mathfrak p)\NK(\mathfrak p)^{-s}\bigr).
\]
Every prime $\mathfrak p\notin S$ has $\NK(\mathfrak p)\ge7$, and $|\nu(\mathfrak p)|\le1$, so for
$\sigma=\Rea s>0$ each factor of $E_\varphi$ is $1+O_\sigma(\NK(\mathfrak p)^{-1-\sigma})$ uniformly
in $\mathfrak p$, with a denominator of absolute value at least $1-7^{-\sigma}$. So the product
converges absolutely, and $E_\varphi(\cdot,\nu)$ is holomorphic on $\Rea s>0$. For $\Rea s>1$,
comparing Euler products gives
\begin{equation}\label{eq:weighted-euler}
 \sum_{(n,S)=1}\frac{\mu(n)\nu(n)\varphi(n)}{\NK(n)^{1+s}}
 =\prod_{\mathfrak p\notin S}\Bigl(1-\frac{\nu(\mathfrak p)}{\NK(\mathfrak p)^{s}}\Bigl(1-\frac1{\NK(\mathfrak p)}\Bigr)\Bigr)
 =\frac1{L_K(s,\nu)P_S(s,\nu)E_\varphi(s,\nu)},
\end{equation}
where $L_K(s,\nu)=L(s,\chi)L(s,\chi\chi_{-3})$. Indeed, a rational prime $p$ that splits in $K$
contributes $(1-\chi(p)p^{-s})^{-2}$ to $L_K(s,\nu)$, an inert one contributes
$(1-\chi(p)^2p^{-2s})^{-1}$, and the ramified prime $3$ contributes $(1-\chi(3)3^{-s})^{-1}$;
these are the Euler factors of $L(s,\chi)L(s,\chi\chi_{-3})$, since $\chi_{-3}(p)=1,-1,0$ in the
three cases.

\begin{proposition}\label{prop:mellin}
Suppose that $A^\varphi(D)\ll_{W,\varepsilon}D^{23/24+\varepsilon}$ for every $\varepsilon>0$ and
every smooth $W:\R\to\C$ that vanishes outside $(1,2)$. Then $L(s,\chi)\ne0$ for
$\Rea s>23/24$, except at the pole $s=1$ when $\chi$ is principal.
\end{proposition}

\begin{proof}
We use the hypothesis to rule out a zero of $L(s,\chi)$ in
$\Rea s>23/24$. Suppose such a zero $\varrho$ exists. Since $L(s,\chi)\ne0$ on $\Rea s\ge1$,
apart from the pole (see \cite[Chapters~13 and~14]{Dav00}), we have $\Rea\varrho<1$.
Choose $0\ne\phi\in C_c^\infty((1,2))$, $\phi\ge0$, and
$W(y)=y^{-\varrho}\phi(y)$.
With the Mellin convention
$\widehat W(s)=\int_0^\infty W(y)y^s\dd y/y$, we have
$\widehat W(\varrho)=\int\phi(y)\dd y/y>0$.
The function
\[
 \mathcal M_W(s)=\int_0^\infty A^\varphi(D)D^{-s}\frac{\dd D}{D}
\]
is holomorphic on $\Rea s>23/24$: the hypothesis
controls the integral at infinity, and
the compact support of $W$ makes $A^\varphi(D)$ vanish for sufficiently
small $D$. Termwise integration for $\Rea s>1$ and \eqref{eq:weighted-euler} give
\[
 \mathcal M_W(s)
 =\widehat W(s)\sum_{(n,S)=1}\frac{\mu(n)\nu(n)\varphi(n)}{\NK(n)^{1+s}}
 =\frac{\widehat W(s)}{L(s,\chi)L(s,\chi\chi_{-3})P_S(s,\nu)E_\varphi(s,\nu)}.
\]
Dirichlet $L$-functions continue meromorphically to $\C$, with at most a simple pole at $s=1$,
which occurs only for principal characters (see \cite[Chapters~8 and~9]{Dav00}); at most one of
$\chi$ and $\chi\chi_{-3}$ is principal. The identity theorem therefore shows that
$L(s,\chi)L(s,\chi\chi_{-3})P_S(s,\nu)E_\varphi(s,\nu)\mathcal M_W(s)=\widehat W(s)$ holds on
$\Rea s>23/24$, $s\ne1$.
Evaluation at $s=\varrho$ gives $0=\widehat W(\varrho)>0$, a contradiction.
\end{proof}

\begin{proof}[Proof of Theorem~\ref{thm:main} from Proposition~\ref{thm:ms}]
Let $\chi$ be a Dirichlet character, let $W$ be smooth and vanish outside $(1,2)$, and let
$\varepsilon>0$. Proposition~\ref{thm:ms}, applied with $I=[1,2]$ and loss $2\varepsilon$, and
Proposition~\ref{prop:extraction} give $A^\varphi(D)\ll_{W,\varepsilon}D^{23/24+\varepsilon}$.
Proposition~\ref{prop:mellin} concludes. For $\zeta$, take $\chi=1$ of modulus $1$; at $s=1$,
where the proposition does not apply, $\zeta(1)$ is a pole, not a zero.
\end{proof}

